\chapter{Metodologia, Métodos e Ferramentas} \label{chp:methodology}
Neste capítulo serão apresentados os princípios metodológicos e ferramentas para tipificar, caracterizar e conduzir a pesquisa. São apresentadas a metodologia aplicada ao estudo em laboratório efetuado, os métodos empregados na condução do estudo e o ferramental utilizado na extração dos dados, para validação das informações elicitadas.

\section{Tipo de Pesquisa}
É chamada pesquisa aplicada aquela que gera conhecimento visando aplicação prática, direcionada para a solução de problemas específicos \cite{sienaMetodologiaPesquisaCientifica2024}. Assim, o presente trabalho busca investigar o desempenho, no que diz respeito à atenção e engajamento dos alunos submentidos a aulas expositivas em ambiente remoto, aplicando e validando o \textit{software} ESPEON, para que,- por meio de relatórios quantitativos - seja possível fundamentar argumentos contrários ao uso da metodologia expositiva, sobretudo em sala de aula remota.

\section{Metodologia Aplicada}
Trata-se de uma pesquisa descritiva explicativa que busca, por meio da observação sistemática, descrever comportamentalmente a população (alunos) submetida à metodologia expositiva em sala de aula virtual. Pretende-se, com o auxílio do ESPEON para coleta e análise de dados, elicitar o desempenho, no que diz respeito ao comprometimento discente, dos alunos submetidos ao cenário proposto, a fim de identificar os desafios que a aula expositiva enfrenta nesse contexto.

É descritiva porque objetiva, primeiramente, caracterizar uma população submetida a um fenômeno (aula expositiva remota), utilizando técnicas padronizadas de coleta de dados (\textit{payload} da \textit{API}), conforme prescrito por \cite{gilComoElaborarProjetos2008a}.

Também é considerada explicativa porque busca, por meio das evidências (relatórios de aula), indicar os principais comportamentos demonstrados pelos alunos, relacionando-os à falta de atenção e engajamento nesse tipo de metodologia. Nesse sentido, reforça-se que:

\begin{customquote}
    A pesquisa explicativa visa explicar a razão dos fatos, por meio da identificação e
    análise das relações de causa e efeito dos fenômenos. Nas ciências sociais, para a coleta de
    dados pode ser utilizado o método observacional ou outros métodos aplicáveis, a depender do
    desenho da pesquisa, dos objetivos e das hipóteses ou suposições formuladas pelo pesquisador.
    \cite{sienaMetodologiaPesquisaCientifica2024}
\end{customquote}

\section{Abordagem da Pesquisa}
No que se refere à abordagem do problema, \cite{ribeiroManagementAccountingAnalysis2015} destaca que uma pesquisa quantitativa é aquela que emprega a quantificação tanto na coleta quanto no tratamento dos dados, por meio de técnicas estatísticas. Neste trabalho, serão utilizadas ferramentas de estatística descritiva, como a média aritmética (medida de tendência central), o desvio-padrão (indicador de dispersão) e a moda (valor de maior frequência), de forma a caracterizar a população estudada. Os cálculos serão processados pelo ESPEON no \textit{backend} e consolidados em gráficos, visando facilitar a interpretação dos resultados e possibilitar análises mais claras.

A escolha da estatística descritiva justifica-se pelo fato de que ela permite sintetizar e interpretar grandes volumes de dados de maneira objetiva, fornecendo subsídios para a comparação entre diferentes comportamentos observados durante as aulas remotas.

% Talvez caiba colocar uma definição desses caras na literatura, e um parágrafo final justificando sua escolha

\section{Método de Pesquisa}
Buscando averiguar a problemática do emprego da metodologia expositiva no ambiente remoto de ensino, parte-se da hipótese de que alunos submetidos a esse tipo de aula, seja no EAD ou no ERE, tendem a apresentar dispersão de foco e baixo engajamento em razão dos inúmero estímulos provenientes de telas, em particular o celular \cite{lopesUsoCelularEm2017}.

Para testar essa hipótese, a pesquisa será conduzida por meio de experimentos controlados com alunos do curso de Sistemas de Informação do CEFET/RJ, campus Maria da Graça. Os participantes serão submetidos às condições mencionadas e terão seu comportamento monitorado pelo sistema ESPEON, mediante consentimento prévio, de forma a possibilitar a coleta de dados para posterior análise estatística.

O método adotado é caracterizado como hipotético-dedutivo, visto que, conforme \cite{ribeiroManagementAccountingAnalysis2015}, fundamenta-se na experimentação para verificar a validade de uma teoria — neste caso, a de que aulas expositivas em ambientes remotos se mostram ineficazes.

\section{Procedimentos e Instrumentos de Coleta de Dados}
A coleta dos dados ocorre de forma passiva durante aulas remotas ministradas no CEFET/RJ, campus Maria da Graça, mediante consentimento dos envolvidos, utilizando o sistema ESPEON. O diagrama abaixo ilustra o processo de coleta, armazenamento e processamento dos dados comportamentais.

\begin{figure}[H]
    \centering
    \begin{minipage}{0.7\textwidth}
        \centering
        \includegraphics[width=.9\textwidth]{espeon/Monitoring.drawio.png}
        \caption{Emissão de \textit{logs}}
        \label{fig:figura1}
    \end{minipage}\hfill
\end{figure}

Ao final das aulas, os dados serão compilados em um relatório, que apresenta numericamente as principais métricas de atenção e engajamento, oferecendo uma avaliação do desempenho da aula ministrada. As métricas avaliadas nesse relatório estão descritas na Tabela \ref{tab:metrics}.

\begin{table}[ht]
    \centering
    \vspace{0.5cm}
    \begin{tabular}{|c|c|}
        \hline
        Métrica                  & Descrição                                                                 \\
        \hline
        Tempo de Inatividade     & Tempo total gasto fora da guia onde ocorre a aula.                        \\
        Tab Swap                 & Número de vezes que o estudante trocou de abas durante a aula.            \\
        Área do Conhecimento     & Segmentação dos dados por área do conhecimento da disciplina ministrada.  \\
        Tempo de Foco            & Tempo máximo de foco contínuo na guia onde ocorre a aula.                 \\
        Aba Silenciada           & Tempo máximo em que a aba da aula permaneceu silenciada.                  \\
        Permissionamento         & Verificação de permissionamento para uso de periféricos de áudio e vídeo. \\
        Streaming de Periféricos & Verificação de uso ativo (streaming) dos periféricos de áudio e vídeo.    \\
        \hline
    \end{tabular}
    \caption{Métricas utilizadas para avaliação de atenção e engajamento}
    \label{tab:metrics}
\end{table}


Essas métricas são construídas pelo processamento de \textit{snapshots} das guias, armazenadas em banco NoSQL, no \textit{backend} Flask. Pode-se observar o formato do \textit{payload} no Script~\ref{script:logs.json} a seguir.

\spaceH{}
\lstinputlisting[
    language=json,
    caption={Exemplo de \textit{log}},
    label={script:logs.json}
]{scripts/mongo.dump.json}

A análise dos dados é realizada com base nos relatórios em PDF emitidos pelo ESPEON (Figura~\ref{fig:report}), que apresentam de forma objetiva métricas de estatística descritiva, bastando apenas sua visualização para interpretação inicial dos resultados.

\begin{figure}[ht]
    \centering
    \begin{minipage}{0.8\textwidth}
        \centering
        \includegraphics[width=.9\textwidth]{espeon/pdf.png}
        \caption{Exemplo de relatório emitido pelo ESPEON}
        \label{fig:report}
    \end{minipage}\hfill
\end{figure}

Adicionalmente, será utilizado o \textit{Power BI} para uma visualização mais intuitiva e interativa dos dados, de modo a compor uma análise de caráter macro sobre o conjunto coletado. A Figura~\ref{fig:dashboard} ilustra o \textit{dashboard} desenvolvido a partir do consumo da \textit{API} de relatórios.

\begin{figure}[ht]
    \centering
    \begin{minipage}{0.999\textwidth}
        \centering
        \includegraphics[width=.9\textwidth]{espeon/relatório.png}
        \caption{\textit{Dashboard} macro de análise}
        \label{fig:dashboard}
    \end{minipage}\hfill
\end{figure}

\section{Ferramentas Utilizadas}\label{sec:tools}
A confecção desta pesquisa foi possibilitada graças a um leque de ferramentas e tecnologias. Esta seção busca elucidar o leitor acerca de todo aparato usado para construção do trabalho, de forma a permitir futura reprodução do estudo.

\subsection{Chrome API}\label{sub:chrome_api}
\subsection{Python Flask}\label{sub:flask}
\subsection{MongoDB e Compass}\label{sub:mongodb}
\subsection{PostgreSQL e DBeaver}\label{sub:postgres}
\subsection{Hosting}\label{sub:hosting}
De forma a fazer os testes em sala de aula, faz-se necessário o \textit{deploy} do ESPEON, tornando-o acessível de maneira centralizada e concisa aos alunos e professores. Considerando as tecnologias utilizadas, \textit{backend} e bancos de dados, foi criado um ambiente na Web, junto às \textit{hostings} Render e AtlasDB, todos usando planos gratuítos de hospedagem.

\subsubsection{Render}
Render é um seriço de hospedagem de aplicações que possuir planos de serviço gratuito para hospedagem de inúmeros tipos de projetos. É especialmente interessante para este projeto porque permite o \textit{deploy} contínuo da aplicação Flask (Figura \ref{fig:render_backend}), mediante conclusão de \textit{Pull Request} e ao mesmo tempo o \textit{hosting} do banco relacional PostgreSQL (Figura \ref{fig:render_postgres}).

\begin{figure}[ht]
    \centering
    \begin{minipage}{0.999\textwidth}
        \centering
        \includegraphics[width=.9\textwidth]{espeon/render_backend.png}
        \caption{Instância \textit{backend}}
        \label{fig:render_backend}
    \end{minipage}\hfill
\end{figure}

\begin{figure}[ht]
    \centering
    \begin{minipage}{0.999\textwidth}
        \centering
        \includegraphics[width=.9\textwidth]{espeon/render_postgres.png}
        \caption{Instância PostgreSQL}
        \label{fig:render_postgres}
    \end{minipage}\hfill
\end{figure}

Então, é um serviço que fornece conectividade remota ao ESPEON, viabilizando os testes em sala de aula propostos neste trabalho, de forma acessível e descomplicada.

\subsubsection{AtlasDB}
AtlasDB é um serviço \textit{multi-cloud} de \textit{hosting} de banco de dados MongoDB. Similar ao Render, destaca-se pela acessibilidade no que diz respeito a custos e implantação, sendo um serviço ideal para \textit{lift and shift} de \textit{databases} locais. Neste projeto, foi utilizada a instância \textit{basic}, que dá direito a 512 MB de armazenamento além de RAM e CPU compartilhadas. Abaixo, na Figura \ref{fig:atlasdb} pode-se observar o \textit{dashboard} do \textit{cluster} ESPEON.

\begin{figure}[ht]
    \centering
    \begin{minipage}{0.999\textwidth}
        \centering
        \includegraphics[width=.9\textwidth]{espeon/atlasdb.png}
        \caption{\textit{Cluster} ESPEON}
        \label{fig:atlasdb}
    \end{minipage}\hfill
\end{figure}

\subsection{LaTeX}\label{sub:latex}

\subsection{Zotero}\label{sub:zotero}
Zotero é um \textit{software} de coleta, organização, anotação, citação e compartilhamento de pesquisas, um gerenciador bibliográfico em suma. A ferramenta foi usada em conjunto com o \textit{plug-in} Zotero Connector de forma a facilitar a organização das fontes bibliográficas usadas na confecção desta pesquisa, contribuindo para sua categorização e alta disponibilidade das fontes, sendo acessível de quaisquer disposivitos \textit{desktop}, \textit{web} e \textit{mobile}. Também possui uma facilidade no \textit{download} das fontes para o formato BibTex, usado neste trabalho. A Figura \ref{fig:zotero} a seguir ilustra o painel geral da aplicação, com as bibliografias segmentadas na barra da esquerda e informações adicionais sobre uma dada fonte ilustradas à direita.

\begin{figure}[H]
    \centering
    \begin{minipage}{0.8\textwidth}
        \centering
        \includegraphics[width=.9\textwidth]{espeon/zotero.png}
        \caption{Aplicação Zotero}
        \label{fig:zotero}
    \end{minipage}\hfill
\end{figure}

\subsection{Git, Github e Actions}\label{sub:git}
Como forma de centralizar e distribuir tanto o \textit{software} ESPEON, quanto esta monografia - mantendo consistência e versionamento de código - foi adotado o uso do Git. Para garantir a disponibilidade, foi escolhido o Github como \textit{hosting} para os projetos. A construção foi dividida em dois repositórios distintos, dada a natureza da operação: um para o \text{software} e outro para a monografia. A gerência desses dois projetos foi feita no Github Projects, o qual prevê estrutura de \textit{cards} similar ao Kanban, com facilidades como documentação de tarefas, acompanhamento ativo do andamento do projeto e organização dos \textit{cards} no geral, conforme ilustrado na Figura \ref{fig:github_board}.

\begin{figure}[H]
    \centering
    \begin{minipage}{0.8\textwidth}
        \centering
        \includegraphics[width=.99\textwidth]{espeon/github.boards.jpeg}
        \caption{\textit{Board} para Monografia}
        \label{fig:github_board}
    \end{minipage}\hfill
\end{figure}

Também foi construída uma \textit{pipeline} de entrega contínua para eventuais \textit{releases} da monografia com o Github Actions, conforme Figura \ref{fig:github_actions} abaixo.

\begin{figure}[H]
    \centering
    \begin{minipage}{0.8\textwidth}
        \centering
        \includegraphics[width=.99\textwidth]{espeon/actions.releases.png}
        \caption{\textit{Releases} da Monografia}
        \label{fig:github_actions}
    \end{minipage}\hfill
\end{figure}

\section{Mecanismo de Análise de Dados}\label{sec:data_analysis}
Visando obter conclusões válidas, deve-se interpretar os dados do experimento \cite{wohlinAnalysisInterpretation2000}. A interpretação dos dados deve ocorrer em três etapas, conforme ilustra a Figura \ref{fig:steps-data}.

\begin{figure}[ht]
    \centering
    \begin{minipage}{0.999\textwidth}
        \centering
        \includegraphics[width=.9\textwidth]{espeon/analise.dados.wohlin.png}
        \caption{Passos para interpretação quantitativa. \textbf{fonte:} \cite{wohlinAnalysisInterpretation2000}}
        \label{fig:steps-data}
    \end{minipage}\hfill
\end{figure}

Segundo \cite{wohlinAnalysisInterpretation2000}, inicialmente os dados são caracterizados por meio da estatística descritiva. Em seguida, realiza-se a exclusão de dados atípicos ou inválidos, de modo a reduzir o \textit{dataset} a pontos confiáveis. Por fim, no terceiro passo, os dados são analisados por meio de testes de hipótese, avaliando-se estatisticamente a conjectura que originou a pesquisa.

Considerando que esta investigação possui caráter quantitativo, optou-se pela análise descritiva dos dados, empregando medidas de tendência central para caracterizar o elemento “aula” analisado.